% Part: model-theory
% Chapter: basics
% Section: partial-iso

\documentclass[../../../include/open-logic-section]{subfiles}

\begin{document}

\olfileid{mod}{bas}{pis}
\section{التشاكلات الجزئية}

\begin{defn}
  لتكن !!{structure}s اثنتان $\Struct{M}$ و$\Struct{N}$. \emph{التشاكل
  الجزئي} من $\Struct{M}$ إلى $\Struct{N}$ دالة جزئية منتهية~$p$ تأخذ
  مدخلاتها من $\Domain M$ وتعيد قيمًا في $\Domain N$، وتحقق شروط التشاكل
  الواردة في \olref[iso]{defn:isomorphism} على مجالها:
  \begin{enumerate}
  \item تكون $p$~!!{injective}؛
  \item لكل !!{constant}~$c$: إذا كانت $p(\Assign{c}{M})$ معرّفة، فإن
    $p(\Assign{c}{M}) = \Assign{c}{N}$؛
  \item لكل !!{predicate}~$P$ ذي $n$ مواضع: إذا كانت $a_1$، \dots،
    $a_n$ في مجال $p$، فإن $\langle a_1, \dots, a_n\rangle \in
    \Assign P M$ إذا وفقط إذا كان
    $\langle p(a_1), \dots, p(a_n) \rangle \in \Assign P N$؛
  \item لكل !!{function}~$f$ ذات $n$ مواضع: إذا كانت $a_1$، \dots،
    $a_n$ في مجال $p$، فإن $p(\Assign f M (a_1, \dots,a_n))
    = \Assign f N (p(a_1), \dots, p(a_n))$.
  \end{enumerate}
  وكون $p$ منتهية يعني أن $\dom{p}$ منتهٍ.
\end{defn}

لاحظ أن الدالة الخالية~$\emptyset$ تشاكل جزئي دائمًا بين أي
!!{structure}s اثنتين.

\begin{defn}\ollabel{defn:partialisom}
  تكون !!{structure}s اثنتان $\Struct{M}$ و$\Struct{N}$
  \emph{متشاكلتين جزئيًا}، ونكتب $\Struct{M} \iso[p] \Struct{N}$، إذا
  وفقط إذا وجدت مجموعة غير خالية~$I$ من التشاكلات الجزئية بين
  $\Struct{M}$ و$\Struct{N}$ تحقق خاصية \emph{الذهاب والإياب}:
  \begin{enumerate}
  \item (\emph{الذهاب}) لكل $p \in I$ و$a \in \Domain M$ يوجد
    $q \in I$ بحيث $p \subseteq q$ ويكون $a$ في مجال $q$؛
  \item (\emph{الإياب}) لكل $p \in I$ و$b \in \Domain N$ يوجد
    $q \in I$ بحيث $p \subseteq q$ ويكون $b$ في مدى $q$.
  \end{enumerate}
\end{defn}

\begin{thm}\ollabel{thm:p-isom1}
  إذا كان $\Struct{M} \iso[p] \Struct{N}$ وكانت $\Struct{M}$ و
  $\Struct{N}$~!!{enumerable}، فإن $\Struct{M} \iso \Struct{N}$.
\end{thm}

\begin{proof}
  لما كانت $\Struct{M}$ و$\Struct{N}$~!!{enumerable}، اكتب
  $\Domain{M} = \{a_0, a_1, \ldots \}$ و
  $\Domain{N} = \{b_0, b_1, \ldots \}$. نبدأ بأي $p_0 \in I$، ونعرّف
  متتالية متزايدة من التشاكلات الجزئية
  $p_0 \subseteq p_1 \subseteq p_2 \subseteq \cdots$ كما يأتي:
  \begin{enumerate}
  \item إذا كان $n+1$ فرديًا، أي إذا كان $n = 2r$، فاستعمل خاصية الذهاب
    لإيجاد $p_{n+1} \in I$ بحيث $p_n \subseteq p_{n+1}$ ويكون $a_r$
    في مجال $p_{n+1}$؛
  \item إذا كان $n+1$ زوجيًا، أي إذا كان $n=2r+1$، فاستعمل خاصية الإياب
    لإيجاد $p_{n+1} \in I$ بحيث $p_n \subseteq p_{n+1}$ ويكون $b_r$
    في مدى $p_{n+1}$.
  \end{enumerate}
إذا وضعنا الآن:
\[
p = \bigcup_{n\ge 0} p_n,
\]
كان $p$ تشاكلًا بين $\Struct{M}$ و$\Struct{N}$.
\end{proof}

\begin{prob}
  بيّن بالتفصيل أن $p$ كما عرّفت في
  \olref[mod][bas][pis]{thm:p-isom1} تشاكل بالفعل.
\end{prob}

\begin{thm}\ollabel{thm:p-isom2}
  لتكن $\Struct{M}$ و$\Struct{N}$~!!{structure}s للـ!!{language} علائقية
  محضة (أي !!a{language} لا تحتوي إلا !!{predicate}s، ولا تحتوي
  !!{function}s أو ثوابت). إذا كان $\Struct{M} \iso[p] \Struct{N}$،
  فإن $\Struct{M} \elemequiv \Struct{N}$ أيضًا.
\end{thm}

\begin{proof}
  يبيّن الاستقراء على الـ!!{formula}s أنه إذا كانت $a_1$، \dots، $a_n$
  و$b_1$، \dots، $b_n$ بحيث يوجد تشاكل جزئي~$p$ يرسل كل $a_i$ إلى
  $b_i$، وكان $s_1(x_i)=a_i$ و$s_2(x_i)=b_i$ (لـ$i=1$، \dots،~$n$)،
  فإن $\Sat{M}{!A}[s_1]$ إذا وفقط إذا كان $\Sat{N}{!A}[s_2]$.
  وتعطي حالة $n=0$ أن $\Struct{M} \elemequiv \Struct{N}$.
\end{proof}

\begin{rem}
إذا وجدت !!{function}s، بقيت النتيجة السابقة صحيحة، لكن يلزم النظر في
التشاكل الذي تستحثه $p$ بين الـ!!{structure} الجزئية من $\Struct{M}$
المولدة بـ$a_1$، \dots، $a_n$ والـ!!{structure} الجزئية من
$\Struct{N}$ المولدة بـ$b_1$، \dots، $b_n$.
\end{rem}

يمكن «تفكيك» النتيجة السابقة إلى مراحل بإقامة صلة بين عدد الكمّامات
المتعشقة في !!a{formula} وعدد المرات التي يمكن فيها توسيع التشاكلات
الجزئية المعنية.

\begin{defn}
  لكل !!{formula}~$!A$، تعرّف \emph{رتبة الكم} لـ$!A$، ويرمز إليها
  بـ$\QuantRank{!A} \in \Nat$، استقرائيًا بأنها أكبر عدد من الكمّامات
  المتعشقة في $!A$. وتكون !!{structure}s اثنتان $\Struct{M}$ و
  $\Struct{N}$ \emph{متكافئتين حتى الرتبة~$n$}، ونكتب
  $\Struct{M} \elemequiv[n] \Struct{N}$، إذا اتفقتا على جميع الجمل ذات
  رتبة الكم الأصغر من أو المساوية لـ$n$.
\end{defn}

\begin{prop}\ollabel{prop:qr-finite}
  لتكن $\Lang{L}$~!!{language} علائقية محضة منتهية، أي !!{language}
  تحتوي عددًا منتهيًا من الـ!!{predicate}s والـ!!{constant}s ولا تحتوي
  !!{function}s. لكل $n,k \in \Nat$ لا يوجد، حتى التكافؤ المنطقي، إلا
  عدد منتهٍ من !!{formula}s الرتبة الأولى في الـ!!{language}~$\Lang{L}$ التي لا تزيد
  رتبة كمها على~$n$ وتقع متغيراتها الحرة ضمن $x_1$، \dots،~$x_k$.
\end{prop}

\begin{proof}
  بالاستقراء على $n$، مع إثبات الدعوى في آن واحد لكل $k$.
\end{proof}

\begin{defn}
  لتكن !!a{structure}~$\Struct{M}$. نرمز بـ$\Domain M^{<\omega}$ إلى
  مجموعة جميع المتتاليات المنتهية على $\Domain{M}$. ونستعمل
  $\mathbf{a}, \mathbf{b}, \mathbf{c}, \ldots$ للدلالة على متتاليات
  منتهية. وإذا كانت $\mathbf{a} \in \Domain{M}^{<\omega}$ و
  $a \in \Domain{M}$، فإن $\mathbf{a}a$ تمثل \emph{وصل}
  $\mathbf{a}$ بـ$a$.
\end{defn}

\begin{defn}
  لتكن !!{structure}s اثنتان $\Struct{M}$ و$\Struct{N}$. نعرّف علائق
  $I_n \subseteq \Domain M^{<\omega} \times \Domain N^{<\omega}$ بين
  متتاليات متساوية الطول، بالاستقراء على $n$ كما يأتي:
   \begin{enumerate}
   \item تكون $I_0(\mathbf{a},\mathbf{b})$ إذا وفقط إذا حققت
     $\mathbf{a}$ و$\mathbf{b}$ الـ!!{formula}s الذرية نفسها في
     $\Struct{M}$ و$\Struct{N}$؛ أي إذا كان طول المتتاليتين~$k$، وكان
     $s_1(x_i)=a_i$ و$s_2(x_i)=b_i$، وكانت $!A$ ذرية وجميع
     !!{variable}s فيها من بين $x_1$، \dots،~$x_k$، فإن
     $\Sat{M}{!A}[s_1]$ إذا وفقط إذا كان $\Sat{N}{!A}[s_2]$.
   \item تكون $I_{n+1}(\mathbf{a},\mathbf{b})$ إذا وفقط إذا كان لكل
     $a\in \Domain M$ عنصر~$b\in \Domain N$ بحيث
     $I_n(\mathbf{a}a,\mathbf{b}b)$، والعكس بالعكس.
   \end{enumerate}
\end{defn}


\begin{defn}
  نكتب $\Struct{M} \approx_n \Struct{N}$ إذا تحققت
  $I_n(\emptyseq,\emptyseq)$ في $\Struct{M}$ و$\Struct{N}$ (حيث
  $\emptyseq$ هي المتتالية الخالية).
\end{defn}

\begin{thm}\ollabel{thm:b-n-f}
  لتكن $\Lang{L}$~!!{language} علائقية محضة. إذا كانت
  $I_n(\mathbf{a},\mathbf{b})$ وكان طول المتتاليتين~$k$، فإن لكل $!A$
  تقع متغيراتها الحرة ضمن $x_1$، \dots،~$x_k$ وتحقق
  $\QuantRank{!A} \le n$، يكون $\Sat{M}{!A}[\mathbf{a}]$ إذا وفقط إذا
  كان $\Sat{N}{!A}[\mathbf{b}]$ (حيث تحقق $\mathbf{a}$ مرة أخرى~$!A$
  إذا حققها أي تعيين~$s$ بحيث $s(x_i)=a_i$). وإذا كانت $\Lang{L}$
  منتهية، صح العكس أيضًا.
\end{thm}

\begin{proof}
  يثبت بسهولة، بالاستقراء على $!A$، أن
  $I_n(\mathbf{a},\mathbf{b})$ تستلزم أن $\mathbf{a}$ و$\mathbf{b}$
  تحققان الـ!!{formula}s نفسها التي لا تزيد رتبة كمها على~$n$ وتقع
  متغيراتها الحرة ضمن قائمة المتغيرات الثابتة المقابلة لإحداثياتهما.
  ولإثبات العكس نستقرئ على $n$، مستعملين
  \olref{prop:qr-finite}، التي تضمن أنه، لكل $n$ ولكل قائمة منتهية ثابتة
  من المتغيرات الحرة، لا يوجد إلا عدد منتهٍ من فئات التكافؤ المنطقي
  للـ!!{formula}s التي لا تزيد رتبة كمها على~$n$.

  عند $n=0$، يقتضي فرض أن $\mathbf{a}$ و$\mathbf{b}$ تحققان
  الـ!!{formula}s الخالية من الكم نفسها أنهما تحققان الصيغ الذرية نفسها،
  ومن ثم تكون $I_0(\mathbf{a},\mathbf{b})$.

  في حالة $n+1$، افترض أن $\mathbf{a}$ و$\mathbf{b}$ تحققان
  الـ!!{formula}s نفسها التي لا تزيد رتبة كمها على $n+1$. ولكي نثبت
  $I_{n+1}(\mathbf{a},\mathbf{b})$، يكفي أن نثبت أنه لكل
  $a \in \Domain M$ يوجد $b \in \Domain N$ بحيث
  $I_n(\mathbf{a}a,\mathbf{b}b)$؛ وبحسب فرض الاستقراء يكفي أيضًا أن
  نثبت أنه لكل $a \in \Domain M$ يوجد $b \in \Domain N$ بحيث تحقق
  $\mathbf{a}a$ و$\mathbf{b}b$ الـ!!{formula}s نفسها التي لا تزيد رتبة
  كمها على~$n$.

  خذ $a \in \Domain M$. اختر ممثلًا واحدًا من كل فئة تكافؤ منطقي
  للـ!!{formula}s~$!B(x,\mathbf{y})$ التي لا تزيد رتبة كمها على~$n$،
  وتقع متغيراتها الحرة ضمن القائمة الثابتة $x,\mathbf{y}$، وتحققها
  $\mathbf{a}a$ في $\Struct{M}$. ولتكن $!T^a_n$ اقتران هذه المجموعة
  المنتهية من الممثلين؛ ومن ثم يجوز عدّها !!{formula} واحدة من الرتبة
  الأولى. يترتب أن $\mathbf{a}$ تحقق
  $\lexists[x][!T^a_n(x,\mathbf{y})]$، ورتبة كم هذه الصيغة لا تزيد على
  $n+1$. وبحسب الفرض، تحقق $\mathbf{b}$ الصيغة نفسها في $\Struct{N}$؛
  ومن ثم يوجد $b \in \Domain N$ بحيث تحقق $\mathbf{b}b$ الصيغة
  $!T^a_n$. وبوجه خاص، تحقق $\mathbf{b}b$ الـ!!{formula}s نفسها التي
  لا تزيد رتبة كمها على~$n$ وتحققها $\mathbf{a}a$؛ إذ تشمل الممثلين
  المختارين ممثل فئة كل صيغة صادقة أو نفي صيغة كاذبة. وبالمثل نبين أنه
  لكل $b \in \Domain N$ يوجد $a\in \Domain M$ بحيث تحقق
  $\mathbf{a}a$ و$\mathbf{b}b$ الـ!!{formula}s نفسها التي لا تزيد رتبة
  كمها على~$n$، وهذا يتم البرهان.
\end{proof}

\begin{cor}\ollabel{cor:b-n-f}
  إذا كانت $\Struct{M}$ و$\Struct{N}$~!!{structure}s علائقية محضة في
  !!{language} منتهية، فإن $\Struct{M} \approx_n\Struct{N}$ إذا وفقط
  إذا كان $\Struct{M} \elemequiv[n] \Struct{N}$. وبوجه خاص،
  $\Struct{M} \elemequiv \Struct{N}$ إذا وفقط إذا كان، لكل $n$،
  $\Struct{M} \approx_n \Struct{N}$.
\end{cor}

\end{document}
